\introduction % Структурный элемент: ВВЕДЕНИЕ

Криптовалюта Bitcoin, созданная в 2008 году Сатоши Накамото~\cite{bitcoin_whitepaper}, продолжает оставаться наиболее капитализированной и широко используемой децентрализованной платёжной системой. В её основе лежит технология распределённого реестра~--- блокчейна~--- представляющего собой публично доступную цепочку блоков транзакций, подтверждённых майнерами посредством алгоритма~\acrshort{pow}. Публичная природа блокчейна обеспечивает полную верифицируемость транзакций, однако участники сети идентифицируются лишь псевдонимными адресами~--- хешами открытых ключей~--- без привязки к реальным личностям. Это создаёт серьёзные вызовы для регуляторов, правоохранительных органов и аналитиков финансовых рисков, заинтересованных в выявлении подозрительной активности: отмывания денег, мошенничества, финансирования незаконной деятельности.

Классификация Bitcoin-адресов по типу деятельности~--- задача, решение которой позволяет частично снять псевдоанонимность участников сети. Если известно, что некоторый адрес принадлежит бирже или майнинговому пулу, аналитик может интерпретировать транзакции с этим адресом и реконструировать финансовые потоки. В отличие от ручного анализа или эвристических методов, основанных на фиксированных правилах, подходы на основе машинного обучения позволяют автоматически выявлять сложные паттерны в пространстве признаков, масштабироваться на большие объёмы данных и адаптироваться к новым видам активности.

В ходе предшествующих научно-исследовательских работ~(НИРС-01 и НИРС-02) были решены задачи систематизации существующих алгоритмов классификации транзакций, сбора Bitcoin-адресов из открытых источников, извлечения 34 транзакционных признаков через~\acrshort{api} Bitcoin-эксплорера и формирования датасета объёмом 8\,115 адресов с разметкой по шести типам кошельков. Выявлена существенная несбалансированность классов~(коэффициент~$\approx$~39,85) и определён набор наиболее информативных признаков.

Настоящая работа является третьим этапом исследовательского цикла и посвящена практическому применению подготовленного датасета: обучению, настройке гиперпараметров и сравнительной оценке моделей классификации Bitcoin-адресов, а также интерпретации полученных результатов.

Актуальность работы обусловлена:
\begin{itemize}
    \item ростом объёма криптовалютных транзакций и потребностью в масштабируемых автоматизированных инструментах мониторинга~--- ручной анализ блокчейна в современных масштабах практически невозможен;
    \item регуляторными требованиями соблюдения процедур \acrshort{aml}/\acrshort{kyc} в криптовалютной сфере, которые обязывают биржи и сервисы идентифицировать контрагентов;
    \item научным интересом к применению ансамблевых методов машинного обучения для анализа высокодимензиональных блокчейн-данных с дисбалансом классов;
    \item необходимостью сравнительного анализа нейросетевых и деревьевидных подходов на табличных Bitcoin-данных.
\end{itemize}

Целью данной научно-исследовательской работы является разработка, обучение и сравнительная оценка ансамблевых и нейросетевых моделей классификации Bitcoin-адресов по типам кошельков на основе датасета, подготовленного в предшествующих НИРС, а также анализ качества и интерпретируемости полученных результатов.

Для достижения цели поставлены следующие задачи:
\begin{itemize}
    \item Подготовка признакового пространства: удаление высококоррелированных признаков, выбор и применение стратегии устранения дисбаланса классов, стандартизация данных;
    \item Обучение шести моделей машинного обучения: MLP, Random Forest, XGBoost, LightGBM, Stacking и kNN;
    \item Настройка гиперпараметров деревьевидных моделей методом случайного поиска с кросс-валидацией;
    \item Оценка качества обученных моделей с использованием метрик взвешенной F1-меры и макро-усреднённого ROC-AUC на отложенной тестовой выборке;
    \item Оптимизация порогов классификации лучшей модели по критерию Юдена~--- для улучшения полноты на миноритарных классах;
    \item Интерпретация предсказаний лучшей модели методом~\acrshort{shap} для выявления наиболее значимых признаков;
    \item Сравнительный анализ всех обученных моделей и формулировка рекомендаций по дальнейшему улучшению.
\end{itemize}
